\hebrewsection{Radar Signal Processing}

Digital Signal Processing is fundamental to modern radar systems. This chapter explores the principles of pulse compression, Doppler processing, and the ambiguity function.

\subsection{Basic Radar Principles}
A radar system transmits a signal $s(t)$ and listens for echoes. The time delay $\Delta t$ between transmission and reception determines the range $R$:
\begin{equation}
    R = \frac{c \Delta t}{2}
\end{equation}
where $c$ is the speed of light.

\subsection{Pulse Compression}
To achieve high range resolution without sacrificing transmit energy, radar systems use long, modulated pulses rather than short, high-power pulses.

\subsubsection{Linear Frequency Modulation (LFM) - Chirp}
The LFM signal, or chirp, increases its frequency linearly over the pulse duration $T$:
\begin{equation}
    s(t) = e^{j\pi K t^2}, \quad -T/2 \leq t \leq T/2
\end{equation}
where $K$ is the chirp rate. At the receiver, a matched filter (pulse compressor) collapses the pulse to a width of approximately $1/B$, where $B=KT$ is the bandwidth.

\subsection{The Ambiguity Function}
The ambiguity function $|\chi(\tau, f_d)|$ represents the output of the matched filter when the received signal has a delay $\tau$ and a Doppler shift $f_d$. It defines the resolution and clutter rejection capabilities of a specific waveform.

\subsection{Doppler Processing and MTI}
Moving Target Indication (MTI) uses the Doppler shift $f_d = \frac{2v}{\lambda}$ to distinguish moving targets from stationary ground clutter. Modern systems use \textbf{Pulse Doppler Radar} with FFT-based processing to resolve targets in both range and velocity.

\begin{pythonbox}{LFM Chirp Generation}
\begin{verbatim}
import numpy as np
import matplotlib.pyplot as plt

T = 10e-6    # 10 microseconds
B = 5e6      # 5 MHz
fs = 50e6    # 50 MHz
t = np.arange(-T/2, T/2, 1/fs)
K = B/T
chirp = np.exp(1j * np.pi * K * t**2)
\end{verbatim}
\end{pythonbox}

\subsection{In-Depth Analysis: Chirp Signals}
Statistical signal processing techniques rely heavily on the estimation of the auto-covariance matrix. Eigen-decomposition of this matrix allows us to separate the signal subspace from the noise subspace, a technique foundational to high-resolution direction-of-arrival (DOA) estimation algorithms like MUSIC and ESPRIT.

Error control coding introduces structured redundancy to the data stream. By mapping the message space into a higher-dimensional codeword space, the receiver can exploit the geometric distance between valid codewords to detect and correct transmission errors, significantly lowering the Bit Error Rate (BER).

\begin{equation}
    \mathbf{R}_{xx} = \mathbb{E}[\mathbf{x}[n]\mathbf{x}^T[n]] = \frac{1}{N} \sum_{k=1}^{N} \mathbf{x}_k \mathbf{x}_k^T
\end{equation}

Multi-rate signal processing involves the alteration of the sampling rate within the digital domain. Decimation and interpolation require stringent anti-aliasing and anti-imaging filters, respectively. Polyphase decompositions of these filters allow for highly efficient implementations that operate at the lower sampling rate.

\begin{pythonbox}{Simulation: Chirp Signals}
\texttt{import numpy as np}
\texttt{import matplotlib.pyplot as plt}
\texttt{}
\texttt{\# Initialize parameters for chirp\_signals}
\texttt{N = 1024}
\texttt{data = np.random.randn(N) + 1j * np.random.randn(N)}
\texttt{signal\_power = np.var(data)}
\texttt{noise = np.random.normal(0, 0.1, N)}
\texttt{processed\_signal = data + noise}
\end{pythonbox}

\vspace{0.5cm}
The mathematical formulation demonstrates the theoretical limits, while the simulation code provides a framework for empirical verification. These tools are critical for evaluating the efficacy of the proposed methodologies under practical constraints.


\subsection{In-Depth Analysis: Pulse Compression}
The integration of machine learning has opened new frontiers in algorithm design. By treating the entire communication or processing chain as a single end-to-end differentiable autoencoder, deep neural networks can learn optimal representations and coding schemes that outperform traditional human-engineered blocks under non-standard channel conditions.

Mathematical modeling of these systems requires an understanding of their asymptotic behavior. By analyzing the limit as the number of samples approaches infinity, we can derive the Cramér-Rao lower bound, which provides the minimum variance for any unbiased estimator. This bound is essential for determining the theoretical efficiency of our algorithms.

\begin{equation}
    H(e^{j\omega}) = \sum_{n=-\infty}^{\infty} h[n] e^{-j\omega n}
\end{equation}

A key component of modern design is the optimization of the objective function. Using stochastic gradient descent or its variants, we can iteratively update the system parameters. The learning rate must be carefully tuned to ensure convergence while avoiding local minima, a phenomenon frequently encountered in highly non-linear parameter spaces.

\begin{pythonbox}{Simulation: Pulse Compression}
\texttt{import numpy as np}
\texttt{import matplotlib.pyplot as plt}
\texttt{}
\texttt{\# Initialize parameters for pulse\_compression}
\texttt{N = 1024}
\texttt{data = np.random.randn(N) + 1j * np.random.randn(N)}
\texttt{signal\_power = np.var(data)}
\texttt{noise = np.random.normal(0, 0.1, N)}
\texttt{processed\_signal = data + noise}
\end{pythonbox}

\vspace{0.5cm}
The mathematical formulation demonstrates the theoretical limits, while the simulation code provides a framework for empirical verification. These tools are critical for evaluating the efficacy of the proposed methodologies under practical constraints.


\subsection{In-Depth Analysis: Doppler Processing}
The spectral efficiency of the system is fundamentally limited by the available bandwidth and the ambient noise floor. According to the Shannon-Hartley theorem, the channel capacity dictates the maximum error-free data rate. Practical systems utilize advanced coding schemes to operate as close to this limit as possible.

In real-world deployments, hardware imperfections such as phase noise, IQ imbalance, and power amplifier non-linearities introduce severe distortions. Digital pre-distortion (DPD) techniques are often employed to digitally compensate for these analog impairments, ensuring the transmitted signal adheres to strict spectral masks.

\begin{equation}
    P_e \leq \sum_{d=d_{free}}^{\infty} a_d Q\left(\sqrt{\frac{2 d R E_b}{N_0}}\right)
\end{equation}

The transient response of the system provides insight into its stability. By observing the pole-zero plot in the complex plane, we can guarantee bounded-input bounded-output (BIBO) stability. For discrete-time systems, this necessitates that all poles reside strictly within the unit circle.

\begin{pythonbox}{Simulation: Doppler Processing}
\texttt{import numpy as np}
\texttt{import matplotlib.pyplot as plt}
\texttt{}
\texttt{\# Initialize parameters for doppler\_processing}
\texttt{N = 1024}
\texttt{data = np.random.randn(N) + 1j * np.random.randn(N)}
\texttt{signal\_power = np.var(data)}
\texttt{noise = np.random.normal(0, 0.1, N)}
\texttt{processed\_signal = data + noise}
\end{pythonbox}

\vspace{0.5cm}
The mathematical formulation demonstrates the theoretical limits, while the simulation code provides a framework for empirical verification. These tools are critical for evaluating the efficacy of the proposed methodologies under practical constraints.


\subsection{In-Depth Analysis: CFAR}
Computational complexity is a major constraint in mobile and embedded applications. The advent of Fast Fourier Transform (FFT) algorithms revolutionized the field by reducing the complexity of discrete convolution from $O(N^2)$ to $O(N \log N)$, making real-time processing feasible.

Statistical signal processing techniques rely heavily on the estimation of the auto-covariance matrix. Eigen-decomposition of this matrix allows us to separate the signal subspace from the noise subspace, a technique foundational to high-resolution direction-of-arrival (DOA) estimation algorithms like MUSIC and ESPRIT.

\begin{equation}
    \mathcal{L}(\theta) = -\frac{1}{M} \sum_{i=1}^{M} \left[ y_i \log(\hat{y}_i) + (1-y_i)\log(1-\hat{y}_i) \right]
\end{equation}

Error control coding introduces structured redundancy to the data stream. By mapping the message space into a higher-dimensional codeword space, the receiver can exploit the geometric distance between valid codewords to detect and correct transmission errors, significantly lowering the Bit Error Rate (BER).

\begin{pythonbox}{Simulation: CFAR}
\texttt{import numpy as np}
\texttt{import matplotlib.pyplot as plt}
\texttt{}
\texttt{\# Initialize parameters for cfar}
\texttt{N = 1024}
\texttt{data = np.random.randn(N) + 1j * np.random.randn(N)}
\texttt{signal\_power = np.var(data)}
\texttt{noise = np.random.normal(0, 0.1, N)}
\texttt{processed\_signal = data + noise}
\end{pythonbox}

\vspace{0.5cm}
The mathematical formulation demonstrates the theoretical limits, while the simulation code provides a framework for empirical verification. These tools are critical for evaluating the efficacy of the proposed methodologies under practical constraints.

